\section{The original sized input and its compatible preliminary geometry}
\label{sec:initial-geometry}

\status{IG1--IG11 construct an original preliminary datum, its integral
norm factors, diagonal spaces, letter identities and exact positive
integer $d_u$. They prove the literal \texttt{exists\_pre\_sized}
conclusion and supply both coordinate/scalar pins after the prescribed
orientation. IG12 constructs the specified \texttt{baseVol} empty-jet
stage. Its chosen-volume definition is distinguished from an assertion
that this number is the geometric covolume. The final paragraphs
reconstruct the recursive-existence implication with its mass-window
and actual recursive-data premises retained.}

The first prime selection below is qualitative. Its purpose is to
construct one actual preliminary input. It gives no prime in a
prescribed short interval, no uniform bound for that prime, and no
large-scale \texttt{ChGood} assertion. In particular its use in the
first eleven paragraphs does not discharge \texttt{LevelPrime}.
All fields, inclusions, frames and norms are those of the displayed
construction. A later coefficient-field extension may preserve them;
it is not identified with an arbitrary previously chosen base.

\paragraph{IG1: partial ray zeta functions from a shifted plane lattice.}
Let $K$ be an imaginary quadratic field, $\mathcal O=\mathcal O_K$,
and $M\geq2$ an integer. The ray group $G$ modulo $M\mathcal O$ is
the group of fractional ideals prime to $M$, modulo principal ideals
generated by elements congruent to $1$ at the primes of $M$.
It is finite. Indeed the ordinary ideal class group is finite by the
short-vector argument in TL14; the kernel of the ray-to-ordinary
class map is a quotient of the finite group
$(\mathcal O/M\mathcal O)^\times$. For this latter assertion, a
principal ideal prime to $M$ is classified, up to a generator
congruent to $1$, by the residue of its generator, modulo residues
of global units. A denominator prime to $M$ can be inverted in
this finite ring. CRT supplies representatives prime to $M$ in
every ordinary ideal class. These observations also identify the
source's finite quotient of the prime-to-$M$ ideal monoid with $G$;
thus the argument applies to any literal \texttt{RayData}, through
its stated quotient isomorphism.

Fix a complex embedding $\sigma:K\hookrightarrow\mathbb C$, and
write $A_0>0$ for the Euclidean covolume of $\sigma(\mathcal O)$.
For a class $c\in G$ choose a nonzero integral ideal $\mathfrak a$,
prime to $M$, in class $c^{-1}$. Since
$\mathfrak a+M\mathcal O=\mathcal O$, choose
$\alpha_0\in\mathfrak a$ with $\alpha_0\equiv1\pmod {M\mathcal O}$.
Set
\[
 \Lambda=\sigma(M\mathfrak a),\qquad z=\sigma(\alpha_0),
 \qquad A=\operatorname{covol}(\Lambda)
       =M^2\mathrm N(\mathfrak a)A_0.                  \tag{IG.1}
\]
The elements of $\mathfrak a$ congruent to $1$ modulo $M$ form
$\alpha_0+M\mathfrak a$: coprimality gives
$\mathfrak a\cap M\mathcal O=M\mathfrak a$.
This coset excludes zero, since $M\mathcal O$ is a proper ideal.
The map
\[
 \alpha\longmapsto(\alpha)\mathfrak a^{-1}
                                                               \tag{IG.2}
\]
maps that coset onto the integral ideals in ray class $c$.
Its fibers are the orbits of
$U_M=\{u\in\mathcal O^\times:u\equiv1\pmod M\}$.
This group is finite, since $|\sigma(u)|=1$ and a discrete lattice
meets the unit circle in finitely many points. Its action is free;
put $w_M=|U_M|>0$. Conversely, if two elements in (IG.2) give the
same ideal their quotient is a unit, and their local congruences
show that this unit belongs to $U_M$. Thus the multiplicity is
exactly $w_M$, not the number of all units unless these coincide.

For $\Re s>1$, counting lattice points in disks gives absolute
convergence and the exact identity
\[
 \zeta_c(s):=\sum_{\substack{\mathfrak b\subset\mathcal O\\
                       (\mathfrak b,M)=1,\ [\mathfrak b]=c}}
                \mathrm N(\mathfrak b)^{-s}
 =\frac{\mathrm N(\mathfrak a)^s}{w_M}
       \sum_{\lambda\in z+\Lambda}|\lambda|^{-2s}.    \tag{IG.3}
\]
For completeness the required counting bound follows by putting
disjoint disks of a fixed sufficiently small radius about lattice
points. All disks with centers in a disk of radius $R$ lie in a
disk of radius $R+O(1)$, so their number is $O(1+R^2)$.
Dyadic annuli then prove convergence in (IG.3), locally uniformly
in $\Re s>1$.

Put $\theta(t)=\sum_{\lambda\in z+\Lambda}e^{-\pi t|\lambda|^2}$.
The Gaussian Fourier transform and Fourier expansion of its
periodization give
\[
 \theta(t)=\frac1{At}
       \sum_{\nu\in\Lambda^*}
          e^{-\pi|\nu|^2/t}e^{2\pi\mathrm i\langle\nu,z\rangle}.
                                                               \tag{IG.4}
\]
Here $\Lambda^*$ is the Euclidean dual lattice for the pairing
appearing in the phase. To justify the formula directly, periodize
$e^{-\pi t|x|^2}$ on a fundamental parallelogram. Its $\nu$th
Fourier coefficient is $A^{-1}t^{-1}e^{-\pi|\nu|^2/t}$, by the
one-dimensional Gaussian integral in each orthogonal coordinate.
Both lattice sums and their derivatives converge uniformly on
compact sets; Fourier summation therefore gives (IG.4).
Because $0\notin z+\Lambda$, $\theta(t)$ decays exponentially as
$t\to\infty$. The nonzero terms of the dual sum give
$\theta(t)-(At)^{-1}=O(t^{-1}e^{-c/t})$ as $t\downarrow0$.
These estimates hold after multiplication by every fixed power
of $t$ and $\log t$.

Consequently
\[
 E_c(s)=\int_1^\infty\theta(t)t^{s-1}\,dt
       +\int_0^1\bigl(\theta(t)-(At)^{-1}\bigr)t^{s-1}\,dt
\]
is entire, by locally uniform differentiated integration.
Mellin integration of each Gaussian in the half-plane of absolute
convergence yields the continuation
\[
 \zeta_c(s)=
  \frac{\mathrm N(\mathfrak a)^s\pi^s}{w_M\Gamma(s)}
          \left(E_c(s)+\frac1{A(s-1)}\right).           \tag{IG.5}
\]
The reciprocal gamma function is entire with simple zeros at the
nonpositive integers. Hence (IG.5) has no pole except a simple
pole at $1$, whose residue is the same positive number for every
class:
\[
          \rho=\frac{\pi}{w_M M^2 A_0}>0.
                                                               \tag{IG.6}
\]
In particular $Z_c(s)=(s-1)\zeta_c(s)$ is entire,
$Z_c(1)=\rho$ and $Z_c(-1)=0$. This proves the actual partial-zeta
continuation assertion, including its common residue and its
negative-integer value. No nonvanishing theorem or prime existence
has been used in this argument.

\paragraph{IG2: nonvanishing at one and prime norms in every ray class.}
Write $h=|G|$. For each complex character $\chi$ of $G$, put
$L_\chi(s)=\sum_c\chi(c)\zeta_c(s)$ initially in $\Re s>1$.
Unique factorization of nonzero ideals gives
\[
 L_\chi(s)=\prod_{\mathfrak p\nmid M}
       (1-\chi([\mathfrak p])\mathrm N(\mathfrak p)^{-s})^{-1}.
                                                               \tag{IG.7}
\]
Absolute convergence follows from (IG.3), summed over the finite
class group. Character orthogonality and (IG.6) show that every
nontrivial $L_\chi$ is entire. The principal $L_1$ has just one
pole, at $1$, with residue $h\rho$, and $L_1(-1)=0$.

Here is the nonvanishing argument with its analytic domain made
explicit. Suppose $L_{\chi_0}(1)=0$ for some $\chi_0\ne1$.
Then $F(s)=\prod_\chi L_\chi(s)$ is entire, since this zero
cancels the only possible pole. If $d_{\mathfrak p}$ is the order
of $[\mathfrak p]$ in $G$, finite character duality gives
\[
 F(s)=\prod_{\mathfrak p\nmid M}
    (1-\mathrm N(\mathfrak p)^{-d_{\mathfrak p}s})^{-h/d_{\mathfrak p}}
      =\sum_{n\ge1}b_n n^{-s}\quad(\Re s>1),
 \qquad b_n\ge0,\quad b_1=1.                          \tag{IG.8}
\]
The coefficient expansion is absolutely convergent in that stated
half-plane. Fix a real $\sigma_0>1$. Termwise differentiation there
is justified because $(\log n)^r\ll_{r,\delta}n^\delta$ for every
$\delta>0$; choose $\delta<\sigma_0-1$. For any real
$\sigma<\sigma_0$, the entire Taylor series of $F$ at $\sigma_0$
has nonnegative terms after inserting these derivatives:
\[
 \begin{split}
 F(\sigma)
 &=\sum_{r\ge0}\frac{(\sigma_0-\sigma)^r}{r!}
           \sum_{n\ge1}b_n(\log n)^r n^{-\sigma_0}\\
 &=\sum_{n\ge1}b_n n^{-\sigma}\ \geq1.               \tag{IG.9}
 \end{split}
\]
Tonelli justifies the interchange; finiteness of the first sum
comes from its equality with the entire Taylor series. At
$\sigma=-1$ this contradicts $F(-1)=0$, forced by the principal
factor. Thus every nontrivial $L_\chi(1)$ is nonzero. This proof
uses the entire continuation obtained under the supposed zero;
it assumes neither an abscissa bound for (IG.8) nor the desired
nonvanishing assertion.

For real $s>1$, expand the logarithm in (IG.7). The terms of
degree at least two have a uniformly bounded total as $s\downarrow1$:
their absolute value is at most a constant times
$\sum_{\mathfrak p}\mathrm N(\mathfrak p)^{-2}<\infty$.
For nonprincipal $\chi$, the full logarithm stays bounded near
$1$. Indeed $L_\chi$ is nonzero on a small disk about $1$ and
has an analytic logarithm there; its difference from the Euler
logarithm on the real interval is a constant multiple of
$2\pi\mathrm i$. For the principal character its logarithm is
$-\log(s-1)+O(1)$ by the positive residue. Orthogonality now gives,
for each $c\in G$,
\[
 \sum_{[\mathfrak p]=c}\mathrm N(\mathfrak p)^{-s}
                   =-\frac1h\log(s-1)+O(1).           \tag{IG.10}
\]
In a quadratic field any non-prime prime-ideal norm is $p^2$.
Their contribution at $s\ge1$ is bounded by
$2\sum_{p}p^{-2}$; there are also only finitely many ideals of
bounded norm. The divergence in (IG.10) therefore leaves, in
every ray class and above any given bound, a prime ideal whose
norm is a rational prime. This is exactly the prime-in-class
conclusion used in \texttt{ExistsNormPrime}.

\paragraph{IG3: integral norm factors and rational progression primes.}
Let $M\ne0$, let $a\in\mathcal O_K$ be coprime to $M$, and let
$N$ be a prescribed bound. For $M\ge2$, apply IG2 to the class
of $(a)$ and choose $\mathfrak p$ with prime norm $p>N$.
Ray equivalence supplies nonzero $x,y\equiv1\pmod M$ with
$x\mathfrak p=y(a)$. The containment
$ya\in x\mathfrak p\subseteq x\mathcal O_K$ gives
$ya=x\pi$ for some $\pi\in\mathcal O_K$.
Cancellation of the nonzero ideal $(x)$ gives
$\mathfrak p=(\pi)$, and
\[
       \pi-a=(y-1)a-(x-1)\pi\in M\mathcal O_K.
                                                               \tag{IG.11}
\]
The field norm is $|\sigma(\pi)|^2>0$, so
$N_{K/\mathbb Q}(\pi)=p$. If $M=1$, use $M=2,a=1$
and discard its extra congruence.

When $K$ is the source host's specified imaginary quadratic
subfield of $E$, map this actual integral element into $\mathcal O_E$.
Conjugation on $E$ restricts to the nontrivial automorphism of $K$;
hence $\pi\bar\pi=p$ in $E$. For a rational integer $a$, write
$\pi=a+M\beta$. Then
\[
 \frac{p-a^2}{M}=a(\beta+\bar\beta)+M\beta\bar\beta
                                                               \tag{IG.12}
\]
is both rational and an algebraic integer, and therefore an integer.
Thus $p\equiv a^2\pmod M$, with the precise integral factor and
the precise congruence required by \texttt{exists\_norm\_prime}.

We also need ordinary rational primes in a reduced residue class;
they can be obtained from the same argument without a new
nonvanishing input. Given $M>0$ and $r$ coprime to $M$, choose a
prime $\ell\equiv3\pmod4$ not dividing $M$. Such primes exist
outside every finite set: $4\prod\ell_i-1$ has an odd number of
prime factors congruent to $3$ modulo $4$ and none of the listed
ones. The quadratic field $K_0=\mathbb Q(\sqrt{-\ell})$ has
discriminant $-\ell$, and is unramified at every prime dividing
$M$, including $2$. At those primes its integral quadratic algebra
is split or unramified. The norm on its units is onto: reduce to
the surjective finite-field norm and lift by the trace correction
in (DR.5)--(DR.6), which works also at $2$; for a split algebra use
$(r,1)$. Choose such local units with norm $r$ and apply CRT to
their residues modulo the prime powers in $M$. This gives
$a\in\mathcal O_{K_0}$ with $N(a)\equiv r\pmod M$ and
$(a,M)=1$. Applying (IG.11) in $K_0$ gives prime norms
$p>N$ with $p\equiv N(a)\equiv r\pmod M$.
Thus all progression-prime selections below follow from IG1--IG3;
no assertion about primes in a shrinking interval has entered.

\paragraph{IG4: an actual cyclic arithmetic host with the required labels.}
The source provides the particularly simple rank record
$R=\texttt{Rt}(0)$: its frame is the cyclic frame of exponent $2$,
\[
       k=8,\qquad m=2,\qquad t=15,\qquad n=4,\qquad f=8.
                                                               \tag{IG.13}
\]
These values satisfy every parity and rank inequality directly:
$r=k+m-1=9$, $r+5\le t$, $3m<t$, $m\le f\le k$ and
$m^2\le k+m(t+1)$. Its map $h$ is the isomorphism
$j\mapsto[5^j]$ from SZ13, and $o=1$, so $o^{64}\le2^e$
with $e=2$.

More generally, for the same cyclic frame at any $e\ge2$, and
given finite bounds $B,N$, put
$P_B=\prod_{0\le i<B}(2i+1)$, with empty product $1$.
It is odd and
$\gcd(8P_B-1,16P_B)=1$. IG3 gives a prime
\[
 d>N+4tr+3,\qquad d\equiv8P_B-1\pmod {16P_B}.
                                                               \tag{IG.14}
\]
Then $d\equiv7\pmod {16}$ and $-d\equiv1\pmod p$ for
every odd prime $p\le B$. The residue $1+6d$ is coprime to
$8d$ and is $3$ modulo $8$. Choose a prime
\[
 q>N+4tr+3,\qquad q\equiv1+6d\pmod {8d}.             \tag{IG.15}
\]
We have $q\equiv1\pmod d$ and $q,d\equiv3\pmod4$.
Quadratic reciprocity therefore gives $(-d/q)=1$:
the signs in $(-1/q)$ and the reciprocity law for $(d/q)$
cancel. These are exactly \texttt{HostHyp}, with $q\ne3$.

Use AH1--AH5 with the trivial prescribed-character group,
$Q_S=2^{e+2}$, and this actual cyclic frame. Concretely the field
is the fixed field of the kernel of
\[
 (\mathbb Z/Q_Sqd\mathbb Z)^\times\longrightarrow
       \{\pm1\}^2\times\mathbb Z/2^e\mathbb Z,
 \quad u\longmapsto\bigl((u/q),(u/d),[u\bmod Q_S]\bigr),
                                                               \tag{IG.16}
\]
where the last class is read through the same $[5]$ generator.
The sign twist of AH3 is trivial here because $[-1]=1$ in the
last quotient. CRT proves surjectivity. AH4 supplies the actual
square roots of $-q,-d$; AH17 supplies all fields of
\texttt{HostDer} for this raw host. Exactness and
$Q_S\le(2^e)^2$ follow from AH5 and $e\ge2$.

Take $\mathfrak c_K=\mathcal O_K$ and the half-different ideal of
AH13. A rational integer $a>1$ is not a unit in $\mathcal O_K$,
as its norm is $a^2>1$; this proves the required primitivity of
$\mathfrak c_K$. Its norm is one, so all its coprimality clauses
hold. The small-splitting assertion is (IG.14); no odd prime
divides $Q_S$. Choose an embedding in the nonempty CM type.
For the initial, subsequently replaced level, the literal choices
$\{31,61,67,71\}$, auxiliary pair $(67,71)$ and distinguished
pair $(31,61)$ satisfy the prime, distinctness, modulo-$15$, and
$Q_S$-coprimality clauses. Set the prescribed finite thick set
to $TT$, the second set to $\{0\}$, the unipotent translate to
$0$, $L_w=2$ and $CT_2=0$. The second-set range holds because
the distinguished cyclic generator has positive order, and its
no-relation assertion has no pair of distinct elements to check.
This constructs an actual \texttt{Host} with the same cyclic
labels and with $q,d>N$; it does not assume an inhabitant of that
record. We use $B=0$ and $TT=\{0\}$ for (IG.13).

\paragraph{IG5: separated auxiliary primes with their integral norm factors.}
Fix the host of IG4 and write $c=Q_Sqd>1$. Choose
$\pi_0\in\mathcal O_E\cap K$ by IG3 with
$\pi_0\bar\pi_0=g_0>N$ prime and $\pi_0\equiv1\pmod c$.
To construct the second prime, define
\[
 P=g_0\prod_{1\le d\le r}(g_0^d-1),\quad
       M=(cP)^{2r+1},\quad
       N'=M^8+1+2M+P.                                 \tag{IG.17}
\]
Both $P$ and $M$ are positive, and $c\mid M$.
Apply IG3 to residue $1+M$ modulo $M^2$ and a bound larger
than $N'+g_0+N$. It gives
\[
 \pi_1\bar\pi_1=g_1,\qquad
 \pi_1\equiv1+M\pmod {M^2},\qquad
 g_1\equiv(1+M)^2\pmod {M^2}.                         \tag{IG.18}
\]
Thus $g_1\ne g_0$, both $g_i>N$, and
$\pi_i\equiv1\pmod c$, $g_i\equiv1\pmod c$.

We check the stronger separation of the actual bad-prime sets.
The source defines
\[
 \mathrm{BB}(r,g)=\{g\}\cup
 \{\ell:\ell\text{ prime},\ \exists\,1\le d\le r,
                   g^d<\ell^{4v_\ell(g^d-1)}\}.       \tag{IG.19}
\]
Every member $\ell$ of $\mathrm{BB}(r,g_0)$ divides $P$.
Put $a=2r+1$; then $\ell^a\mid M$ and $2d<\ell^a$ for
$1\le d\le r$. From the bound and congruence in (IG.18) write
$g_1=1+2M+M^2b$, $b\ge0$, and set
$x=M(2+Mb)$ and $S=\sum_{i=0}^{d-1}(1+x)^i$.
The geometric-sum identity and reduction modulo $M$ give
\[
 g_1^d-1=M(2+Mb)S,
       \qquad (2+Mb)S\equiv2d\pmod M.                 \tag{IG.20}
\]
Since $0<2d<\ell^a$, $\ell^a$ does not divide $(2+Mb)S$.
Consequently
\[
 \ell^{4v_\ell(g_1^d-1)}
 \le (\ell^a)^4(\ell^{v_\ell(M)})^4
 \le M^8\le g_1\le g_1^d.                            \tag{IG.21}
\]
This excludes its exceptional-prime clause in (IG.19), and
$g_1>P$ excludes $\ell=g_1$. Hence the two \texttt{BB} sets
are disjoint, with the literal fourth-power valuation threshold.

Finally use IG3's rational progression conclusion to choose
$p_\delta>g_1$ and $p'_\delta>p_\delta$, both
$1\pmod {15c}$. Replace the level by
$\{g_0,g_1,p_\delta,p'_\delta\}$ and use these ordered pairs
as the auxiliary and distinguished primes. Congruence to $1$
modulo $c$ proves coprimality with $Q_S$; all four primes are
distinct as required. The field, Artin map, half-different,
thick set and rank labels are unchanged. This proves the actual
releveling conclusion and all the integral factor congruences.
For IG4's fixed rank choose $N=2fk+1$, which gives
$2fk<g_i-1$.

\paragraph{IG6: two distinct coded columns and a common spread bound.}
Choose the two constant Boolean words on $\{0,\ldots,k-3\}$,
one true and one false; $k\ge m+3$ makes this set nonempty.
Put $\Delta=1$, $h=t+3$, $u_F=4(k-2)$, $u_L=-1$ and $b_0=1$.
For either code $c$, define
\[
 \begin{split}
 v_i&=4(k-2-i) &&(0\le i<k-1),\\
 u_0&=4(k-2)+\tfrac12,\qquad u_{k-1}=-\tfrac12,\\
 u_{i+1}&=v_i-\tfrac12-c_i &&(0\le i<k-2),\\
 b_j&=1+2jh,\qquad e_j=b_j+h &&(0\le j<m).
 \end{split}                                                   \tag{IG.22}
\]
These formulas are independent of the CM embedding. Every
$u_i$ is a half-integer; every $v_i,b_j,e_j$ is integral.
At an interior position, $v_i-u_{i+1}$ is $1/2$ or $3/2$,
and $u_{i+1}-v_{i+1}$ is $7/2$ or $5/2$.
The endpoints have gaps $1/2$. Also
$u_{k-1}=-1/2>-b_0=-1$ and
$e_j-b_j=b_{j+1}-e_j=h$, which satisfies the source's
gap requirement because $2h\ge t+3$.
These are all inequality and integrality fields of
\texttt{HCParam}. The common $v,b,e$, first and last values
give the five \texttt{PreCh} compatibility equations.
The two columns are distinct because their value at $u_1$
differs by one. Their letters obey exactly
\[
       v_i-v_{i+1}=4\Delta,\qquad
       u_{i+1}=v_i-\tfrac12-\Delta c_i.                \tag{IG.23}
\]
For each $j$, the exponent $e_j=1+(2j+1)(t+3)$ is constant
over the CM type. AC2--AC3 show that its product on every
root of unity is one. This is the literal \texttt{MUQ}
condition in \texttt{PreX.e\_even}; that field's name does
not impose an additional parity assumption on $e_j$.

There are finitely many triples consisting of a word, an index
$0\le j\le m$, and a CM embedding. Let $S$ be the maximum of
the real spreads of their actual weights $\lambda_j$.
Choose an integer $a\ge1$ sufficiently large that
\[
                 S\le2a,\qquad 2fN_m+1<4a.           \tag{IG.24}
\]
For example an integer strictly above
$\max(1,S/2,(2fN_m+1)/4)$ works. This proves all spread
inequalities simultaneously, with one $a$ fixed before the
type prime is chosen. It is a bound on the actual finite weight
family, and does not replace its integral Weyl lattices.

\paragraph{IG7: a common type prime with an exact initial window.}
Apply IG3 with modulus
\[
                  M_T=Q_Sqd(g_0g_1)^2
\]
and residue $1$, choosing a prime $T$ larger than
$4a+p_\delta+p'_\delta$. Its integral factor $\pi_T\in K$
satisfies $\pi_T\bar\pi_T=T$ and
$\pi_T-1\in M_T\mathcal O_E$.
Use this same $T,\pi_T$ for both words. It is distinct from
both distinguished primes; it is $1$ modulo $Q_Sqd$ by (IG.12).
Set initially
\[
 d_u=1,\qquad D_*=4T^{fk},\qquad\Gamma_P=2.
                                                               \tag{IG.25}
\]
All are positive and $\Gamma_P>1$. Since $T>0$ and $fk>0$,
the real-power rules give exactly
\[
     (D_*/d_u)^{1/(fk)}\Gamma_P^{-2/(fk)}=T.
                                                               \tag{IG.26}
\]
Thus the entire two-sided window is
$T\le T\le(1+1/(fk))T$.
Together with primality, field membership, norm, congruence
and avoidance, these are all fifteen fields of \texttt{LevCh}.
Extend the existing level by the image of this constant type-prime
map; the host field and its labels remain unchanged.
The scalar in \texttt{Pre.ofCh} is the literal quotient
$c(u)=\pi_T(u_0)/\pi_T(u)=1$. In a nonconstant variant the
same quotient obeys
$c(u)\overline{c(u)}T_u=T_{u_0}$ by cancellation of its
nonzero denominator. No new rational norm equation is assumed.

\paragraph{IG8: construct the spaces and every field of the original PreX.}
Use the actual CM extension and involution of IG4, and the norm
factor $p_z=\pi_T$ of IG7. The coordinate construction of SE15
gives $X_j=E^{N_j}$ and $X'=E^{\operatorname{rank}Y_m}$ with
positive diagonal forms and zero-extension inclusions. In $Y_j$
put a positive sign on its first
$N_j+\lfloor(m-1-j)/2\rfloor$ coordinates and a negative
sign on the remainder. The hyperbolic pairs are exactly
$e_i+f_i,(e_i-f_i)/2$ in the matched positive/negative slots.
Their pairings are $0,0,1$, respectively. An extra negative
coordinate, when present, contains $p_z$ and has norm $-T$.
At the top the extra positive coordinate contains $p_z$ and
has norm $T$. All the asserted orthogonalities follow because
their coordinate supports are disjoint. Take
$W_j$ with signature $(k+j-1,m-j)$ and every $V_j$ positive;
their coordinate lattices are $\mathcal O_E$ in each slot.

There is also a finite-adelic isometry field in the actual
\texttt{Spaces} record. It must be supplied, not inferred from
an archimedean signature. AH11 gives the literal
$\mathrm{cmUnram}$ of this extension. DR1 proves that at every
finite base-field prime the integral CM algebra is split or
unramified, and its unit norm is onto. Choose a unit $y_v$
with $\bar y_vy_v=-1$ there. These choices are integral units
at every finite place and therefore form a finite idele $y$,
with $\bar yy=-1$. To change one diagonal sign multiply that
coordinate by $y$; to keep a sign multiply by $1$.
The inverse multiplier is $-\bar y$, and entrywise calculation
gives $\bar g^{\mathsf T}J_sg=J_t$ for any two of these sign
patterns. This supplies the literal
\texttt{baseSpace\_isometric\_graded} on the same finite CM
algebra. Together with SE15 it verifies all the remaining
space fields and \texttt{DiagSpaces}, including
$V\mathrm{neg}=\mathrm{false}$ for every index.

Construct $P=\texttt{Pre.ofCh}$ with these spaces and the two
columns of IG6. Its base word is the first code. Now define
$X:\texttt{PreX}(P)$ using precisely $\pi_0,\pi_1$ from IG5,
$\pi_T$ from IG7, and $a$ from (IG.24).
The six auxiliary/type factor membership, norm and congruence
fields are already proved. The distinguished primes are
$1$ modulo the host conductor $Q_Sqd$, so the cyclotomic
Frobenius formula AH9 gives complete splitting in $E$.
Thus each has exactly $[E:\mathbb Q]$ primes above it, which
is the exact \texttt{pde\_split} cardinality clause.
The field \texttt{ccK\_top} is IG4; the \texttt{MUQ} clause
and injectivity of the weight family are IG6.
Finally $2fk<g_i-1$ is IG5, $4a<T_u$ is IG7,
and the two remaining size/spread inequalities are (IG.24).
These exhaust the fields of \texttt{PreX}.
The word set has cardinal two, $TT=\{0\}$ is nonempty, and
(IG.23) gives the injective two-letter family. This reconstructs
\texttt{exists\_pre\_two\_words} from its underlying arithmetic
and coordinate geometry.

\paragraph{IG9: the exact mass-index expression is a positive integer.}
We next prove the integrality needed before resetting $d_u$.
For a prime $g$ let $U_k(\mathbb F_g)$ be the upper
unitriangular group and write
\[
 i_g=[\mathrm{GL}_k(\mathbb F_g):U_k(\mathbb F_g)]>0,
 \qquad |\mathrm{GL}_k(\mathbb F_g)|
                         =g^{k(k-1)/2}i_g.             \tag{IG.27}
\]
This is exactly the real quantity \texttt{idxColLoc}, since
there are $k(k-1)/2$ independently chosen strict-upper entries.

Let $L$ be a full positive definite Hermitian lattice of rank
$k$ over $\mathcal O_E$, and let $\Gamma$ be its unitary
automorphism group. It is finite by CR4's bounded-image argument.
Assume $g>[E:\mathbb Q]k+1$ and that a prime of $E$ over $g$
has residue field $\mathbb F_g$. Reduction of $L$ there has
dimension $k$: localize its torsion-free finite module at this
DVR, choose a basis, and reduce that basis. Reduction embeds
$\Gamma$ into $\mathrm{GL}_k(\mathbb F_g)$.
Indeed any finite subgroup of the principal congruence kernel
is a $g$-group. The proof is the binomial induction in CR4:
an element of order prime to $g$ congruent to $1$ is congruent
to $1$ at every depth and hence is $1$.
The exponent of this finite-order element is a unit in the local
DVR, so the induction works with arbitrary ramification; it does
not identify the rational prime with a local uniformizer.
But an element of order $g$ would have a primitive $g$th root
as an eigenvalue on the rational space of dimension
$[E:\mathbb Q]k$, forcing $g-1\le[E:\mathbb Q]k$.
Cauchy's theorem therefore excludes a nontrivial kernel,
and also proves $g\nmid|\Gamma|$.
Lagrange and (IG.27) now give $|\Gamma|\mid i_g$.

There are finitely many classes in the genus of $L$.
One direct justification uses the trace form: it is a positive
definite integral quadratic form of fixed rank and determinant.
Successive-minimum reduction, with its positive integral minimum,
gives an integral basis with bounded Gram entries; there are
finitely many such Gram matrices. For each one the matrices
giving multiplication by a fixed integral basis of $\mathcal O_E$
are bounded in the trace norm (their archimedean operator norms
are fixed), so only finitely many integral matrices occur.
The Hermitian form is recovered from the trace pairings with
these multiplications by the nondegeneracy of the field trace.
Thus the Hermitian, not merely quadratic, isometry classes are
finite. This is the finite-genus input to the source's mass.
Consequently
\[
       \operatorname{mass}(L)i_g
          =\sum_{[L']\in\operatorname{genus}(L)}
                         \frac{i_g}{|\operatorname{Aut}(L')|}
               \in\mathbb Z_{>0}.                    \tag{IG.28}
\]
The genus contains $L$ itself, so the sum is nonempty. This
uses the finite sum defining the definite mass, not an analytic
Siegel mass formula.

For a strictly decreasing half-integral tuple $u_1,\ldots,u_k$,
the source's dimension factor is
\[
              D(u)=\frac{\prod_{a<b}(u_a-u_b)}
                                   {\prod_{j=1}^{k-1}j!}
                       \in\mathbb Z_{>0}.             \tag{IG.29}
\]
Here is an integral proof. Subtract the common half to get
integers $z_i$, and use the integer-valued polynomials
$\binom{X}{j}$, $0\le j<k$. Their evaluation matrix at the
integers $z_i$ is integral, and its determinant is their
Vandermonde divided by $\prod_{j=0}^{k-1}j!$.
The sign is fixed by the decreasing order, so (IG.29) is
positive. This argument includes negative $z_i$, for which
the binomial polynomial still has integral values.

Apply (IG.28) to the standard positive form and standard lattice
over the actual $E$. Congruence $g_0\equiv1\pmod {Q_Sqd}$
gives a residue-degree-one prime by AH9, and
$[E:\mathbb Q]=2f$ makes IG5's size bound sufficient.
Thus, for every supplied HC column,
\[
 \begin{split}
 d^\mathrm{true}(u)
 &=\operatorname{mass}(\mathcal O_E^k)
                       i_{g_0}^{\,f}i_{g_1}^{\,f}
                        \prod_\tau D(u_\tau)\\
 &=(\operatorname{mass}(\mathcal O_E^k)i_{g_0})
                       i_{g_0}^{\,f-1}i_{g_1}^{\,f}
                        \prod_\tau D(u_\tau)
                         \in\mathbb Z_{>0}.           \tag{IG.30}
 \end{split}
\]
The first line is literally \texttt{duStarOf}. Since $f\ge1$,
every factor in the second line is a positive integer. This
proves \texttt{duStarOf\_nat} without replacing a dimension
by an integral-lattice equality.

\paragraph{IG10: the original sized preliminary input and its exact window.}
Take the actual two-word $P,X$ of IG8 and keep its base word
$u_0$. Let $d=d^\mathrm{true}(u_0)\in\mathbb Z_{>0}$ from
IG9. Restrict the word type to a singleton, keeping its old
base column, type prime, integral norm factor, host, spaces,
coordinate frames, and the same $a$. Retain $\Gamma_P$ and
set
\[
      d_u=d,\qquad D_* =T_{u_0}^{fk}\Gamma_P^2d.
                                                               \tag{IG.31}
\]
All quantities in the real-power calculation are positive.
Consequently
\[
 (D_*/d)^{1/(fk)}\Gamma_P^{-2/(fk)}
  =T_{u_0}\Gamma_P^{2/(fk)}\Gamma_P^{-2/(fk)}=T_{u_0}.
                                                               \tag{IG.32}
\]
Thus both original window inequalities hold, including their
unchanged upper factor $1+1/(fk)$. The inherited norm
quotient scalar and its identities are unchanged at the
base word. Restrict its letter function too; a function from
a singleton is injective, and the two equations (IG.23)
still hold. The fields of \texttt{PreX} independent of words
are copied; its type factors, $T$ bound and spreads are the
old fields at $u_0$. Injectivity of the weight family is now
automatic.

We have therefore constructed, with no supplied preliminary
input, an actual $P$ satisfying
\[
 \mathrm{HasLetters}(P),\quad
       \mathrm{PreX}(P)\ne\varnothing,\quad
       \mathrm{DiagSpaces}(P.\mathrm{sp}),\quad
       d_u=d^\mathrm{true}(u)\quad\text{for every word }u.
                                                               \tag{IG.33}
\]
This is the literal \texttt{JetsInhab.exists\_pre\_sized}
conclusion. The source intentionally reduces the two-word
family to a singleton here; (IG.33) does not claim two
distinct words with a shared newly normalized $D_*$.
The latter would require simultaneous, externally prescribed
windows and is a different statement.

\paragraph{IG11: orientation and both pins on this same preliminary datum.}
Let a prescribed \texttt{PiBar} choose a prime
$\mathfrak p_i$ of the imaginary quadratic subfield over
each $g_i$. The factor $\pi_i$ of IG5 and its conjugate
generate maximal ideals of norm $g_i$ in that subfield.
The chosen prime contains one of them. Replace $\pi_i$ by
itself or its conjugate so that
$(\bar\pi_i)=\mathfrak p_i$ there. For every prime $w$ of
$E$,
\[
     \bar\pi_i\in w\quad\Longleftrightarrow\quad
                         w|_K=\mathfrak p_i.           \tag{IG.34}
\]
This follows by contraction and maximality; it is the exact
orientation condition. The norm is unchanged and conjugation
preserves the congruence to one modulo the rational integer
$Q_Sqd$. Only these auxiliary factor choices change. All
type factors, weights, windows and coordinate frames remain
the ones in IG10.

The source scalar $c(u)=\pi_T(u_0)/\pi_T(u)$ satisfies
$c(u)\pi_T(u)=\pi_T(u_0)$, since its denominator has nonzero
prime norm. The coordinate inclusion $X'\hookrightarrow X_m$
takes the $a$th standard frame vector to the vector with the
same coordinate index; evaluate every coordinate to check the
identity. These are exactly \texttt{Pins.cof} and
\texttt{Pins.iX}. Thus IG10 supplies the previously retained
original input to SE15--SE16, and TL16 applies to any prescribed
orientation. No arbitrary rational frame or previously fixed
type-prime factor is identified with this choice.

The existence of a full Hecke-character frame is a separate
question: orientation here holds for every supplied
\texttt{PiBar} (in particular one coming from such a frame).
It does not manufacture the frame's $\mu_j,\eta,\chi_\Psi,
\delta,\delta'$ by renaming their defining record. The next-source
record identifies their actual supplier if that additional
unconditional join is required.

\paragraph{IG12: the exact chosen-volume empty-jet stage.}
Fix an actual base $B$, its columns $C$, and a compact-dual Haar
system $\mathrm{cd}$, constructed for example by SE11.
For a Hermitian space $W$, the source first defines
$\mathrm{VolIs}(W,x)$ using all Haar measures normalized by
the finite lattice stabilizer and the specified archimedean
compact-dual measures, and all their rational fundamental
domains. Its actual definition of \texttt{trueVol} is
\[
 v(W)=
 \begin{cases}
 \text{a chosen }x\text{ with }x>0\text{ and }\mathrm{VolIs}(W,x),
                    &\text{if such an }x\text{ exists},\\
 1,                 &\text{otherwise}.
 \end{cases}                                                   \tag{IG.35}
\]
Classical choice in the first branch retains its two stated
properties; the second branch is positive as well. Therefore
$v(W)>0$ for every $W$, without an assumption of finite
geometric volume. If a positive \texttt{VolIs} witness has
separately been proved, the first branch applies and
$\mathrm{VolIs}(W,v(W))$ follows. The positivity argument alone
does not give this extra conclusion.

The literal $\mathrm{baseVol}(B,\mathrm{cd},W)$ is (IG.35) for
$B$'s actual CM embedding family and the supplied same
$\mathrm{cd}$. Apply the model construction of SE12 with this
positive function and with both jet index types empty.
Its two adjacent-rank antecedents are then empty universal
statements. This gives an actual $J:\texttt{GJStage}(B,C)$
whose fields satisfy
\[
  J.\mathrm{cd}=\mathrm{cd},\qquad
  J.\mathrm{vol}=\mathrm{baseVol}(B,J.\mathrm{cd}),\qquad
  J.\mathrm{jet}=\mathrm{noJet}(B),
                                                               \tag{IG.36}
\]
and every prescribed \texttt{gat} has finite classes.
The volume equality uses the same Haar datum after substitution;
it is not merely equality with a constant positive model.
This proves the exact \texttt{GJInhab.exists\_noJet} endpoint.
It also gives \texttt{nonempty\_stage} after choosing SE11's
Haar system. For \texttt{exists\_oneJet} the identical argument
retains its two nonempty adjacent-class inputs and its specified
prime/rank situation. The chosen-volume construction supplies
neither those classes nor the geometric \texttt{VolTieG} theorem.

\paragraph{IG13: the exact recursive sizes before imposing their window.}
For an arbitrary supplied $\texttt{PreCh}$, write $d_P(u)$ for
the expression in (IG.30), and let $D_P$ be the source's
\texttt{DsP}:
\[
 D_P=\operatorname{mass}(\mathcal O_E^r)
        \bigl(j_{g_0}^{\,f}j_{g_1}^{\,f}\bigr)
        \prod_\tau d_\infty(\mathrm{hc}(u_0,\tau)),
 \qquad r=k+m-1.                                      \tag{IG.37}
\]
Here $j_g$ is the literal row-level index and
$d_\infty=\texttt{trueDsig}$ is the defined formal-degree
polynomial. The identity matrix belongs to the finite row-level
group, so its order and the order of $\mathrm{GL}_r(\mathbb F_g)$
are both positive. Their quotient $j_g$ is positive. The
concatenation
\[
       (v_0,\ldots,v_{k-2},-b_0,\ldots,-b_{m-1})
\]
is strictly decreasing: the interlacing gives decreasing $v$,
the two $b,e$ gaps give increasing $b$, and
$v_{k-2}>u_{k-1}>-b_0$ gives the intervening strict gap.
The source's $\lambda_S$ is exactly this ordered concatenation,
through the finite-index equality $r=(k-1)+m$. Its
Vandermonde divided by the positive factorial product is
$d_\infty$. Thus each factor in (IG.37) is positive, by the
finite-genus argument of IG9, and
\[
 D_P>0,\qquad
 (2/3)^f\operatorname{mass}(\mathcal O_E^r)
       j_{g_0}^{\,f}j_{g_1}^{\,f}
       \prod_\tau D(v_\tau,-b_\tau)\ \le D_P.           \tag{IG.38}
\]
In fact the product following $(2/3)^f$ equals $D_P$ exactly.
This proves \texttt{DsP\_pos\_and\_lower}; it uses the defined
formal-degree polynomial, and asserts no realization of a
discrete series or its analytic formal-degree theorem.

The integer assertion for $d_P(u)$ is also uniform over the
large-scale inputs occurring in the recursive arrow. Indeed
the already established rank and exponential-versus-polynomial
bounds give, at one threshold depending only on
$0<\epsilon<1/52$,
\[
       f\le x^3,\qquad b\ge x^4,\qquad x\ge10.
\]
If the supplied \texttt{ChOK} has $g_i\ge(kb)^8$, then
$g_i\ge kb\ge10kx^3>2fk+1$, using $k\ge1$.
Its arithmetic clauses give $g_i\equiv1\pmod {Q_Sqd}$.
Apply IG9 separately to every actual word. This gives positive
integers $d_u$ with $d_u=d_P(u)$, at a threshold independent of
the word, $c_D,c_h$, coefficient field or later level choice.
This is the literal \texttt{exists\_du\_nat} conclusion.
The separate fixed-rank existence construction of IG10 is not
being used to replace these supplied large-scale inputs.

\paragraph{IG14: the recursive-existence join and its exact remaining inputs.}
We reconstruct the quantifier and same-object step in
\texttt{recExP\_of\_recAtG}. Its original proof invokes the
following quantitative theorem \texttt{window\_lower}; in
this component that theorem remains an explicit input:
for fixed $\epsilon,C,c_D\ge0$ in its stated range, once $c_h$
and then $x$ exceed its thresholds, every supplied
\texttt{ChOK} and the exact choices $d_u=d_P(u)$,
$D_*=D_P$, $\Gamma=\mathrm{GamS}(C,s,R)$ satisfy
\[
 m\mathcal L\ \le\
  \log\left((D_*/d_u)^{1/(fk)}\Gamma^{-2/(fk)}\right)
                       \quad\text{for every }u.       \tag{IG.39}
\]
The source derives this from \texttt{SiegelMassRatioCM},
the local-index ratio and an averaged Vandermonde lower bound.
IG13 supplies its positivity and size-normalization inputs.
Neither IG7's freely chosen window nor (IG.38) proves (IG.39).
In particular, our reconstruction of the arrow below is
conditional on this remaining quantitative theorem; it is
not a claim to have removed it from the original source proof.

Suppose the supplied $\mathrm{RecAtG}$ gives, over every
sufficiently large good context $K$, an actual $D,Y$ with
\[
 \mathrm{Over}(D,K.P,K.Lv),\quad
 \mathrm{RecOK},\quad\mathrm{LedOK},\quad\mathrm{TWAt},
 \quad\mathrm{FromCoreAt}(Q',K,D,Y),\quad
 \mathrm{TupUsedAt},\quad\mathrm{TwAllAt},\quad
 \mathrm{MtStableAt}.                                  \tag{IG.40}
\]
Also suppose the supplied $\mathrm{EvP}$ asserts $Q_O(\kappa,\beta.D,\beta.O)$
for every sufficiently large actual $\beta:\mathrm{BPt}(K)$
whose core satisfies $Q'$, under $\kappa.\mathrm{Pos}$.
We explain all threshold choices, since changing their order
would change the theorem.

Take the minimum of the two positive upper bounds for $\epsilon$.
For fixed $\epsilon,C_8,C_5,c_6,C_b$, take the supremum of their
two transported lower-bound records; its components are maxima,
so a standard $C$ large for this record is large for each.
For this $C$, take
\[
 c_1=\max(c_1^{\mathrm{rec}},c_1^{Q},0).
\]
After $c_D\ge c_1$ is fixed, take the maximum of the two
$c_h$ thresholds and the threshold of (IG.39).
After $c_h$ is fixed, take the maximum of the four $x$ thresholds
for (IG.40), the $Q_O$ conclusion, (IG.39), and IG13's integral
$d_P$. If $\kappa.\mathrm{Pos}$ is false, use threshold $0$
for the second of these four: the asserted implication is
then vacuous. This conditional choice occurs only after all
components of $\kappa$ have been fixed.

Now let $s,P,a,b,c$ satisfy the resulting good-scale antecedents.
Choose $d_u=d_P(u)\in\mathbb Z_{>0}$ by IG13 and choose
$D_*=D_P>0$. Equation (IG.39) supplies the source's logarithmic
window with the literal constant $c_s=1$.
These choices precede the next universal quantifier.
For \emph{every} subsequently supplied
$Lv:\mathrm{LevCh}(P)$ with
$Lv.d=d$, $Lv.D_*=D_P$ and $Lv.\Gamma=\mathrm{GamS}$, form
\[
 K=(s,C,c_D,c_h,P,q_1,a,b,c,\mathrm{Std},\mathrm{ChOK},
                                        Lv,d_P,D_P,\mathrm{GamS}).
                                                               \tag{IG.41}
\]
The last three entries mean the exact defining equalities of
the context; $q_1$ is the one in the supplied good datum.
Thus this is the actual context over the same $P,Lv$, with
no newly selected prime or word relabeling. Apply (IG.40)
to obtain its $D,Y$.

Fix any permitted $\mathrm{LedPart}$ of $Y$, with the supplied
scale, weight, numerical-value and canonical-instance equations.
The $\mathrm{FromCoreAt}$ field is a witness $X,f,T,\mathfrak c$
and an equality of dependent pairs
\[
             (D,\mathrm{layObj}(Y))
                     =(X.\mathrm{toD7},\mathfrak c.O).
                                                               \tag{IG.42}
\]
Transport first along this literal equality. The fixed
$\mathrm{LedPart}$, its scale/weight/value/canonical equations,
the good datum, $\mathrm{RecOK}$ and $\mathrm{LedOK}$, and the
two used/twist fields of (IG.40) now form exactly a
$\beta:\mathrm{BPt}(K)$. Its $Q'$ condition is the same core
condition in (IG.42). The supplied $\mathrm{EvP}$ therefore
gives $Q_O$ for its actual point and object whenever
$\kappa.\mathrm{Pos}$ holds.

The target also requires $\mathrm{QD}$ after the point's
\texttt{setDiv} operation. Here
$\mathrm{QD}=\mathrm{TupUsedAt}\wedge\mathrm{TwAllAt}
\wedge\mathrm{MtStableAt}$. Changing only the ramified
divisor lattices and their height parameters leaves the
tuple labels, their use, the coefficient module
$\widetilde M$, and its integral scalar action unchanged.
For the twist clause, keep the original conjugacy witness,
the integers $N_0$, and the same finite selection sets.
The definitions of the automorphic record and the twist
epsilon used in every $\mathrm{TwGoodOne}$ are unchanged;
the paired local membership statements transport by their
specified identity maps. Thus the same witnesses prove
the new $\mathrm{TwGoodOne}$ and $\mathrm{QD}$.
This is the actual \texttt{qd\_setDiv} implication, not a
fresh arithmetic choice of twist data.

We have proved the complete \texttt{RecExP} quantifier and
transport conclusion under (IG.39), (IG.40), and the stated
$\mathrm{EvP}$ premise, including its $\mathrm{withQD}$
conjunct. The original source arrow has only
$\mathrm{SiegelMassRatioCM}$ as its extra numerical premise
because it separately invokes \texttt{window\_lower}.
That quantitative deduction, and the supply of actual
recursive cores, remain explicit continuation targets here.

For precision the upstream \texttt{LayOfCore.lay} does more
than join bare auxiliary stages. It starts with an actual
$\mathrm{V4.Core}$ and $\mathrm{VolumeEqMass}_2$ on the same
$X.\mathrm{toD7}$. Its core specifies the true scalar
degrees, the real metric/volume ties, actual local row and
column laws, the common comparison maps, the ramified local
lattice bounds, and the selected arithmetic character data.
The two size equalities of the resulting $\mathrm{Lay}$
come from \texttt{sizes\_eq\_count}; they are not forced by
positive volume or by a bare finite-dimensional model.
Its $\mathrm{RecOK}$ has nine literal duties: the specified
$\Delta$, $\eta_F$, untwisted row vanishing, finite twist
support, positive $D_\Gamma$, fullness of the ramified
integral lattice, tame distinctness, the actual Minkowski
$D_\Gamma$ choice, and the two-sided type-prime local bounds.
Its $\mathrm{LedOK}$ additionally compares the
actual integral $\mathrm{NP}$ lattices for every admissible
canonical instance. These are retained at their actual
producer ranges. IG10--IG12 do not silently supply them.
